\documentclass{article}
\usepackage{amsmath}
\usepackage{amssymb}
\usepackage[margin=1in]{geometry}
\usepackage{enumitem}
\usepackage{graphicx}
\newcommand{\sectionbreak}{\clearpage}
\newcommand{\qed}{$\hfill \square$}


\title{Mastery Workbook 5}
\author{Joe Ricks}
\date{\today}

\begin{document}

\maketitle

\textbf{\Huge I have neither given nor received unauthorized assistance.}

\sectionbreak

\textbf{\#1 (10 pts)}

\bigbreak

Look up the RSA algorithm and write 3 specific questions or observations you have about how it works.

\bigbreak

\begin{enumerate}
    \item

    \bigbreak

    \item

    \bigbreak

    \item

\end{enumerate}

\sectionbreak

\textbf{\#2 (30 pts)}

\bigbreak

Summarize essential definitions and concepts from Sec. 4.1 - 4.3.

\bigbreak



\sectionbreak

% Additional page for Question 2 summary
\textbf{\#2 (continued)}

\bigbreak



\sectionbreak

\textbf{\#3 (10 pts)}

\bigbreak

Read ``The Sieve of Eratosthenes'' on page 259. Notice it uses THM 2, when it says ``note that the composite integers not exceeding 100 must have a prime factor not exceeding 10.''

\bigbreak

\begin{enumerate}[label=\Alph*.]
    \item Consider Theorem 2 page 258 - show inductively that it is true with some examples.

    \bigbreak



    \bigbreak

    \item Re-Prove Theorem 2 page 258 in your own words, using our proof structure for this proof by contradiction.

    \bigbreak



\end{enumerate}

\sectionbreak

\textbf{\#4 (10 pts)}

\bigbreak

Recall that we can think of modulo in terms of scoops and cups. The Euclidean Algorithm can also apply to scoops and cups and solves the question: ``Given 2 measuring scoops, of $a$ cups and $b$ cups, what is the smallest measure in cups that can be made?''

\bigbreak

You must bake a cake and need to measure exactly 1 cup. Can you make the cake given:

\bigbreak

\begin{enumerate}[label=\alph*.]
    \item 2 scoops measuring 17 and 88?

    \bigbreak



    \bigbreak

    \item 2 scoops measuring 35 and 88?

    \bigbreak



    \bigbreak

    \item 2 scoops measuring 37 and 53?

    \bigbreak



\end{enumerate}

\sectionbreak

\textbf{\#5 (10 pts)}

\bigbreak

Following the Square and Mod video, create a ``Square and Mod Chart'' for powers of 5 mod 17.

\bigbreak

Show some EVIDENCE OF THE PROCESS you are using to get the values below.

\bigbreak

\begin{align*}
5^0 \bmod 17 &= \\[6pt]
5^1 \bmod 17 &= \\[6pt]
5^2 \bmod 17 &= \\[6pt]
5^4 \bmod 17 &= \\[6pt]
5^8 \bmod 17 &= \\[6pt]
5^{16} \bmod 17 &= \\[6pt]
5^{32} \bmod 17 &= \\[6pt]
5^{64} \bmod 17 &= \\[6pt]
5^{128} \bmod 17 &= \\[6pt]
5^{256} \bmod 17 &=
\end{align*}

\sectionbreak

\textbf{\#6 (10 pts)}

\bigbreak

Use the chart and the Square 'n Mod method from the video to find:

\bigbreak

\begin{itemize}
    \item $5^{270} \bmod 17$

    \bigbreak



    \bigbreak

    \item $5^{186} \bmod 17$

    \bigbreak



    \bigbreak

    \item $5^{411} \bmod 17$

    \bigbreak



\end{itemize}

\sectionbreak

\textbf{\#7 (10 pts)}

\bigbreak

Choose an FME algorithm to code:
\begin{itemize}
    \item Book Version
    \item Sriram's Version
    \item Your own based on Square 'n Mod
\end{itemize}

\bigbreak

Briefly, explain why you chose which algorithm you did.

\bigbreak



\bigbreak

Include a screenshot of your COMMENTED code below:

\bigbreak

% \includegraphics[width=\textwidth]{fme_code.png}

\sectionbreak

\textbf{\#8 (10 pts)}

\bigbreak

To motivate our use of Fast Modular Exponentiation FME, let's do a comparison of two ways we can calculate $a^n \bmod m$.

\bigbreak

Write a new function called NOT\_FME that returns the result of $a^n \bmod m$, using just the standard exponent calculation and the built in Python Mod function \%.

\bigbreak

Then use the Python function pow(), which calculates the same result using FME.

\bigbreak

Do your own informal investigation to see if FME has a greater impact on the runtime when $n$ is getting large or when $a$ is getting large.

\bigbreak

\textbf{Summary of Results:}

\bigbreak



\end{document}
